\section{Незалежний відтворюваний експеримент для спектрально-зорового та структурного методу}
\label{sec:ch4_ind_svs}

\subsection{Межі експерименту та його зв'язок із ретроспективною перевіркою}
\label{subsec:ch4_ind_svs_scope}

Для усунення частини обмежень ретроспективної перевірки проведено окремий відтворюваний експеримент із пооб'єктним реєстром походження, груповим поділом даних, неперервними діагностичними оцінками та первинними часовими вимірюваннями. Його початкове значення генератора випадкових чисел дорівнювало $20260717$. Усі наведені далі числові результати стосуються лише цього фінального циклу. Початкове значення $20260716$, використане під час пілотного налагодження, до підсумкового оцінювання не входило.

Новий цикл не є повторним обчисленням раніше розглянутого прогону на $110$ об'єктах. Для попереднього прогону збереглася лише агрегована матриця $TP=50$, $TN=48$, $FP=7$, $FN=5$, тоді як у новому експерименті сформовано інші об'єкти, інший протокол поділу та пооб'єктні оцінки. Тому показники двох циклів не об'єднуються, а їх відмінність не тлумачиться як зростання або зниження ефективності одного й того самого випробування.

Предметом незалежної перевірки була здатність структурних і спектрально-зорових ознак розрізняти чисті об'єкти та об'єкти з контрольовано внесеними змінами. Позитивна мітка означала лише наявність такої зміни. Вона не означала наявності виконуваного шкідливого коду, не відтворювала реальне шкідливе програмне забезпечення й не підтверджувала його виявлення. Експеримент також не охоплював метод поведінкового узгодження ознак різного походження, описаний у підрозділі~3.3, і не містив подієво-поведінкових журналів.

\subsection{Набір даних і протокол групового поділу}
\label{subsec:ch4_ind_svs_protocol}

Контрольований набір сформовано детерміновано з $200$ первинних груп: $160$ груп PNG та $40$ груп MP4. Кожна група містила чистий об'єкт і створений з нього змінений об'єкт, тому загальна кількість зразків становила $400$. Розмір PNG дорівнював $256\times256$ пікселів. Відео MP4 мало розмір кадру $160\times120$ пікселів, складалося з $12$ кадрів і кодувалося з частотою $12$ кадрів/с. Базові кадри були синтетичними: їх утворено з гармонічних складників, шуму та геометричних фігур; у відео положення фігур змінювалося між кадрами.

Одиницею поділу була первинна група, а не окремий файл. Спочатку групи стратифіковано за форматом і типом запланованої зміни та розподілено у співвідношенні $60:20:20$, і лише після цього створено чистий і змінений об'єкти кожної групи. Отже, два похідні від одного базового вмісту об'єкти не могли потрапити до різних частин. Склад поділу наведено в табл.~\ref{tab:ch4_ind_svs_dataset}.

\begin{table}[htbp]
\caption{Склад незалежного контрольованого набору та групового поділу}
\label{tab:ch4_ind_svs_dataset}
\centering
\small
\begin{tabular}{lrrrrrr}
\toprule
Формат & Групи & Чисті & Змінені & Навчальна & Валідаційна & Тестова \\
\midrule
PNG & 160 & 160 & 160 & 192 & 64 & 64 \\
MP4 & 40 & 40 & 40 & 48 & 16 & 16 \\
\midrule
Разом & 200 & 200 & 200 & 240 & 80 & 80 \\
\bottomrule
\end{tabular}
\end{table}

У навчальній, валідаційній і тестовій частинах містилося відповідно $120$, $40$ і $40$ первинних груп. Кожна частина була збалансованою за класом, оскільки кожну групу представляла одна чиста й одна змінена версія. У тестовій частині було $40$ негативних і $40$ позитивних об'єктів; з них $64$ об'єкти мали формат PNG, а $16$ --- MP4.

Контрольовані зміни належали до чотирьох типів. Кожному типу відповідало $50$ первинних груп у повному наборі: $30$ навчальних, $10$ валідаційних і $10$ тестових. Способи внесення змін узагальнено в табл.~\ref{tab:ch4_ind_svs_changes}.

\begin{table}[htbp]
\caption{Типи контрольовано внесених змін}
\label{tab:ch4_ind_svs_changes}
\centering
\small
\begin{tabularx}{\textwidth}{p{0.18\textwidth}Y Y}
\toprule
Тип & PNG & MP4 \\
\midrule
\texttt{overlay} & Додавання $2048$ байтів після логічного завершення PNG & Додавання $4096$ байтів після наявного контейнера \\
\texttt{private\_block} & Додавання приватного допоміжного блока \texttt{stEg} обсягом $2048$ байтів перед \texttt{IEND} із коректною контрольною сумою & Додавання контейнерного блока \texttt{free} з корисним навантаженням $4096$ байтів \\
\texttt{spectral} & Малоамплітудна зміна синього каналу: шахове відхилення на $\pm2$ для $18$~\% пікселів і зміна молодшого біта для $12$~\% значень каналу & Покадрове шахове відхилення синього каналу на $\pm(2\ldots4)$ для $16$~\% пікселів \\
\texttt{combined} & Спектральна зміна разом із хвостовим навантаженням $2048$ байтів & Спектральна зміна разом із хвостовим навантаженням $4096$ байтів \\
\bottomrule
\end{tabularx}
\end{table}

Для кожного об'єкта збережено ідентифікатор зразка, ідентифікатор первинної групи, частину поділу, формат, клас, тип зміни, розмір файла, початкове значення генератора й контрольну суму SHA-256. Окремі машинозчитувані записи містили всі виділені ознаки, неперервні оцінки трьох діагностичних конфігурацій, пороги, рішення та часові вимірювання. Контроль цілісності фінальних результатів підтвердив $400$ рядків реєстру, $200$ груп, декодування $320$ PNG і $80$ MP4 та $240$ рядків тестових прогнозів, тобто по $80$ для кожної з трьох конфігурацій.

\subsection{Ознаки та діагностична модель-перевірка}
\label{subsec:ch4_ind_svs_probe}

Структурне подання складалося з дев'яти ознак. До нього входили індикатори помилки сигнатури, синтаксичного розбору та цілісності; відношення хвостових байтів і приватного корисного навантаження до розміру файла; частка невідомих елементів; масштабована кількість структурних елементів; індикатор невідповідності розміру; індикатор відеоконтексту. Для PNG перевіряли сигнатуру, послідовність блоків, CRC, наявність \texttt{IEND}, приватні й невідомі блоки та байти після логічного завершення. Для MP4 послідовно розбирали контейнери верхнього рівня, їх розміри, типи, блоки \texttt{free} і залишок після останнього коректного блока.

Спектрально-зорове подання складалося із шістнадцяти ознак: середньої ентропії; середнього модуля та стандартного відхилення лапласіана; середнього й максимального значення енергії деталізації Гаара; середньої частки високочастотної енергії дискретного косинусного перетворення у блоках $8\times8$; зміщення й частоти переходів молодших бітів; відповідних лапласіанних, гаарових, косинусних і бітових характеристик синього каналу; середнього та максимального міжкадрового відхилення; міжканального залишку; індикатора відеоконтексту. Для PNG часові ознаки дорівнювали нулю, оскільки об'єкт містив один кадр. Спільне подання містило $24$ ознаки: дев'ять структурних і п'ятнадцять додаткових спектрально-зорових, оскільки спільний індикатор відеоконтексту не дублювався.

Для перевірки інформаційної придатності ознак використано однакову для всіх конфігурацій стандартизовану логістичну регресію з $L_2$-регуляризацією $\lambda=0{,}05$. Ця модель виконувала роль відтворюваної діагностичної проби, а не нового складника методу. Середні значення, масштаби та ваги визначали лише за $240$ навчальними об'єктами. Порогове значення вибирали лише за $80$ валідаційними об'єктами за максимумом збалансованої правильності; у разі рівності послідовно враховували $F_1$, близькість порога до $0{,}5$ і менше значення порога. Тестові $80$ об'єктів не використовували ані для стандартизації, ані для навчання, ані для добору порога.

Було навчено три проби: за структурними ознаками, за спектрально-зоровими ознаками та за їх спільним поданням. Валідаційні пороги дорівнювали відповідно $0{,}328738$, $0{,}481738$ і $0{,}456468$. Під час тестування пороги залишали незмінними. Такий експеримент перевіряє роздільну та спільну діагностичну інформативність явно обчислених ознак, але не є навчанням або перевіркою зорового перетворювача (Vision Transformer, ViT). Отже, наведені результати не можна приписувати архітектурі ViT.

Для кількісного оцінювання використано матрицю рішень, правильність (Accuracy), збалансовану правильність (Balanced Accuracy), точність позитивного рішення (Precision), повноту (Recall), специфічність, $F_1$, частоти хибнопозитивних і хибнонегативних рішень та площу під ROC-кривою (AUC). Невизначеність оцінювали процентильним груповим бутстрепом із $2000$ повторень: у кожному повторенні з поверненням вибирали $40$ тестових груп і включали обидва об'єкти кожної вибраної групи. Межами $95$~\% довірчого інтервалу були $2{,}5$-й і $97{,}5$-й процентилі. Пороги під час повторного вибирання не змінювали.

Парні відмінності правильності рішень спільної конфігурації та кожної окремої конфігурації оцінено двобічним точним критерієм Мак-Немара на тих самих $80$ тестових об'єктах. Критерій використовував лише дискордантні пари рішень. Оскільки він застосований на рівні об'єктів, а чистий і змінений об'єкти пов'язані спільним походженням, це порівняння доповнює групові довірчі інтервали, але не замінює зовнішньої перевірки на незалежних природних даних.

\subsection{Результати на відкладеній тестовій частині}
\label{subsec:ch4_ind_svs_results}

Результати трьох конфігурацій на незмінній тестовій частині наведено в табл.~\ref{tab:ch4_ind_svs_metrics}. Спільне подання дало $TN=40$, $FP=0$, $FN=2$, $TP=38$, тобто $\mathrm{Acc}=0{,}9750$, $\mathrm{Rec}=0{,}9500$ і $\mathrm{AUC}=0{,}9869$. Структурна проба пропустила шість змінених об'єктів і сформувала одне хибнопозитивне рішення. Спектрально-зорова проба без структурних ознак пропустила $17$ змінених об'єктів і сформувала сім хибнопозитивних рішень.

\begin{table}[htbp]
\caption{Тестові результати діагностичних конфігурацій незалежного експерименту}
\label{tab:ch4_ind_svs_metrics}
\centering
\small
\resizebox{\textwidth}{!}{%
\begin{tabular}{lrrrrrrrrrrrrr}
\toprule
Конфігурація & $TN$ & $FP$ & $FN$ & $TP$ & Acc & BAcc & Prec & Rec & Spec & $F_1$ & FPR & FNR & AUC \\
\midrule
Структурна & 39 & 1 & 6 & 34 & 0,9125 & 0,9125 & 0,9714 & 0,8500 & 0,9750 & 0,9067 & 0,0250 & 0,1500 & 0,9337 \\
Спектрально-зорова & 33 & 7 & 17 & 23 & 0,7000 & 0,7000 & 0,7667 & 0,5750 & 0,8250 & 0,6571 & 0,1750 & 0,4250 & 0,7500 \\
Спільна & 40 & 0 & 2 & 38 & 0,9750 & 0,9750 & 1,0000 & 0,9500 & 1,0000 & 0,9744 & 0,0000 & 0,0500 & 0,9869 \\
\bottomrule
\end{tabular}}
\end{table}

Порівняно зі структурною пробою спільна конфігурація збільшила правильність на $6{,}25$ відсоткового пункту, повноту --- на $10{,}00$ відсоткового пункту та AUC --- на $0{,}0531$. Порівняно зі спектрально-зоровою пробою приріст правильності становив $27{,}50$ відсоткового пункту, повноти --- $37{,}50$ відсоткового пункту, а AUC --- $0{,}2369$. Ці різниці характеризують саме три повторно навчені логістичні проби в межах сформованого контрольованого набору. Вони не є оцінкою вилучення складника з готової промислової моделі.

Групові довірчі інтервали спільної конфігурації подано в табл.~\ref{tab:ch4_ind_svs_ci}. Для правильності інтервал становив $[0{,}9375;1{,}0000]$, для повноти --- $[0{,}8750;1{,}0000]$, для AUC --- $[0{,}9613;1{,}0000]$. Нульова спостережувана частота хибнопозитивних рішень у тесті спричинила вироджений бутстреп-інтервал $[0;0]$ для FPR. Він означає відсутність хибнопозитивних рішень серед наявних $40$ тестових негативних об'єктів, але не доводить нульової частоти таких рішень у генеральній сукупності.

\begin{table}[htbp]
\caption{Групові 95~\% довірчі інтервали спільної конфігурації}
\label{tab:ch4_ind_svs_ci}
\centering
\small
\begin{tabular}{lrr}
\toprule
Показник & Точкова оцінка & 95~\% довірчий інтервал \\
\midrule
Accuracy & 0,9750 & $[0{,}9375; 1{,}0000]$ \\
Balanced Accuracy & 0,9750 & $[0{,}9375; 1{,}0000]$ \\
Precision & 1,0000 & $[1{,}0000; 1{,}0000]$ \\
Recall & 0,9500 & $[0{,}8750; 1{,}0000]$ \\
Specificity & 1,0000 & $[1{,}0000; 1{,}0000]$ \\
$F_1$ & 0,9744 & $[0{,}9333; 1{,}0000]$ \\
FPR & 0,0000 & $[0{,}0000; 0{,}0000]$ \\
AUC & 0,9869 & $[0{,}9613; 1{,}0000]$ \\
\bottomrule
\end{tabular}
\end{table}

Розподіл результатів спільної конфігурації за типом контрольованої зміни та форматом наведено в табл.~\ref{tab:ch4_ind_svs_disaggregation}. Для \texttt{overlay}, \texttt{private\_block} і \texttt{combined} виявлено всі $10$ позитивних тестових об'єктів кожного типу. Для \texttt{spectral} виявлено $8$ із $10$. За форматом усі $64$ тестові PNG класифіковано правильно, тоді як серед $16$ MP4 отримано два хибнонегативні рішення. Оскільки загалом було лише два пропуски, обидва вони одночасно належали до типу \texttt{spectral} і формату MP4.

\begin{table}[htbp]
\caption{Результати спільної конфігурації за типом зміни та форматом}
\label{tab:ch4_ind_svs_disaggregation}
\centering
\small
\begin{tabular}{llrrrrrr}
\toprule
Розріз & Категорія & $n$ & $TN$ & $FP$ & $FN$ & $TP$ & Показник \\
\midrule
Тип зміни & \texttt{overlay} & 10 & -- & -- & 0 & 10 & $\mathrm{Rec}=1{,}000$ \\
Тип зміни & \texttt{private\_block} & 10 & -- & -- & 0 & 10 & $\mathrm{Rec}=1{,}000$ \\
Тип зміни & \texttt{spectral} & 10 & -- & -- & 2 & 8 & $\mathrm{Rec}=0{,}800$ \\
Тип зміни & \texttt{combined} & 10 & -- & -- & 0 & 10 & $\mathrm{Rec}=1{,}000$ \\
\midrule
Формат & PNG & 64 & 32 & 0 & 0 & 32 & $\mathrm{Acc}=1{,}000$ \\
Формат & MP4 & 16 & 8 & 0 & 2 & 6 & $\mathrm{Acc}=0{,}875$ \\
\bottomrule
\end{tabular}
\end{table}

Точний критерій Мак-Немара показав неоднакову доказову силу двох порівнянь (табл.~\ref{tab:ch4_ind_svs_mcnemar}). У порівнянні зі структурною пробою спільна конфігурація виправила п'ять рішень і не погіршила жодного, однак за п'яти дискордантних пар двобічне значення $p=0{,}0625$ не перетнуло рівень $0{,}05$. Отже, на цій тестовій частині числову перевагу над структурною пробою зафіксовано, але статистично значущої парної відмінності за прийнятим рівнем не встановлено. У порівнянні зі спектрально-зоровою пробою спільна конфігурація виправила $23$ рішення й погіршила одне; значення $p=2{,}98\cdot10^{-6}$ указує на парну відмінність у межах цього контрольованого тесту.

\begin{table}[htbp]
\caption{Парне порівняння правильності рішень за точним критерієм Мак-Немара}
\label{tab:ch4_ind_svs_mcnemar}
\centering
\small
\begin{tabularx}{\textwidth}{Y Y Y c c}
\toprule
Порівняння & Спільна правильна, інша помилкова & Спільна помилкова, інша правильна & Дискордантні & $p$ \\
\midrule
Спільна проти структурної & 5 & 0 & 5 & 0,0625 \\
Спільна проти спектрально-зорової & 23 & 1 & 24 & $2{,}98\cdot10^{-6}$ \\
\bottomrule
\end{tabularx}
\end{table}

\subsection{Час формування ознак}
\label{subsec:ch4_ind_svs_timing}

Час структурного та спектрально-зорового оброблення вимірювали окремо для кожного з $400$ об'єктів високоточним монотонним лічильником. Сумарний час дорівнював сумі двох охоплених операцій і не містив часу генерації даних, навчання, добору порога, реєстрації рішення, роботи черги або передавання мережею. Результати наведено в табл.~\ref{tab:ch4_ind_svs_timing}.

\begin{table}[htbp]
\caption{Пооб'єктний час формування ознак у незалежному експерименті}
\label{tab:ch4_ind_svs_timing}
\centering
\small
\begin{tabular}{llrrrr}
\toprule
Формат & Операція & $n$ & Середнє, мс & Медіана, мс & 95-й процентиль, мс \\
\midrule
PNG & Структурна & 320 & 0,32 & 0,30 & 0,41 \\
PNG & Спектрально-зорова & 320 & 53,57 & 52,86 & 60,01 \\
PNG & Сумарна & 320 & 53,89 & 53,17 & 60,35 \\
\midrule
MP4 & Структурна & 80 & 0,21 & 0,20 & 0,26 \\
MP4 & Спектрально-зорова & 80 & 198,55 & 197,75 & 212,55 \\
MP4 & Сумарна & 80 & 198,76 & 197,95 & 212,74 \\
\midrule
Разом & Структурна & 400 & 0,30 & 0,30 & 0,40 \\
Разом & Спектрально-зорова & 400 & 82,57 & 53,70 & 203,85 \\
Разом & Сумарна & 400 & 82,86 & 54,01 & 204,06 \\
\bottomrule
\end{tabular}
\end{table}

Для PNG середній сумарний час становив $53{,}89$~мс, а для MP4 --- $198{,}76$~мс. В обох форматах переважну частину виміряного часу займало спектрально-зорове оброблення. Для всіх об'єктів разом середній сумарний час дорівнював $82{,}86$~мс, однак це зважене значення залежить від співвідношення $320$ PNG до $80$ MP4 і не повинно переноситися на потік з іншим складом форматів. Повний сценарій, до якого входили формування набору, виділення ознак, навчання трьох проб, оцінювання та створення результатів, виконувався $61{,}62$~с.

Фінальний запуск виконано у Windows~11 з Python~3.13.3, NumPy~2.3.0, OpenCV~4.10.0, Matplotlib~3.9.2 і вбудованою підтримкою FFmpeg. Апаратну конфігурацію у фінальному паспорті середовища не зафіксовано. Тому часові значення характеризують конкретну програмну реалізацію та цей запуск, але не є апаратно незалежною оцінкою швидкодії або наскрізного часу промислової системи.

\subsection{Обмеження та предметне тлумачення результатів}
\label{subsec:ch4_ind_svs_limits}

Новий експеримент усунув змішування похідних одного первинного вмісту між навчальною, валідаційною і тестовою частинами та забезпечив пооб'єктні оцінки, групові довірчі інтервали, парні порівняння й часовий журнал. У межах сформованого набору результати підтримують висновок, що спільне структурне та спектрально-зорове подання містить більше діагностичної інформації, ніж суто спектрально-зорове подання. Порівняно зі структурним поданням спостережуване поліпшення не досягло рівня $p<0{,}05$ за точним критерієм Мак-Немара, тому його слід перевірити на більшій кількості незалежних груп.

Водночас усі первинні зображення та відео були синтетичними, а типи змін --- відомими й алгоритмічно сформованими. Набір охоплював лише PNG і короткі MP4; природні зображення, довгі відео, інші кодеки, JPEG, WebP, повторне кодування соціальними платформами та невідомі способи модифікації не перевірялися. Тест містив лише $40$ первинних груп, зокрема $8$ груп MP4, тому оцінки за форматом і типом зміни мають ширшу невизначеність, ніж загальний показник.

Використана стандартизована логістична регресія була контрольною моделлю-перевіркою явно заданих ознак. Експеримент не перевіряв ViT або іншу глибоку навчувану зорову архітектуру, не оцінював метод підрозділу~3.3, не містив реальних подієво-поведінкових послідовностей і не встановлював здатності виявляти реальне шкідливе програмне забезпечення. Одержані $97{,}50$~\% правильності та AUC $0{,}9869$ характеризують лише розпізнавання чотирьох видів контрольованої зміни у відкладеній частині цього синтетичного набору.

Отже, незалежний цикл надає відтворювану емпіричну перевірку структурних і спектрально-зорових ознак методу в обмежених контрольованих умовах. Його результати доповнюють, але не замінюють ретроспективний аналіз окремого прогону на $110$ об'єктах. Наступна зовнішня перевірка має використовувати природні носії з незалежних джерел, більшу кількість груп і форматів, різні рівні та способи зміни, окрему апаратну специфікацію й незмінний заздалегідь зареєстрований протокол.
